% =====================================================================
%  related_work_table_nodeps.tex
%  Landscape of accelerators for structured sparse GEMV.
%
%  ZERO-DEPENDENCY VARIANT of related_work_table.tex.
%
%  This version needs NO new packages in the ntua-thesis template:
%    - \checkmark comes from amssymb, which ntua-thesis.cls loads in
%      BOTH font branches (cls lines 687 and 692), so it is present
%      whether you compile with "modern" or "classic".
%    - \textcolor comes from xcolor, pulled in by the unconditional
%      \usepackage{tikz} at cls line 1238.
%    - Rules use plain \hline instead of booktabs.
%    - The 3-line header is built by hand instead of with multirow.
%
%  Use related_work_table.tex instead if you install pifont, booktabs
%  and multirow -- it looks better (proper booktabs rule weights).
%
%  USAGE
%    % =====================================================================
%  related_work_table_nodeps.tex
%  Landscape of accelerators for structured sparse GEMV.
%
%  ZERO-DEPENDENCY VARIANT of related_work_table.tex.
%
%  This version needs NO new packages in the ntua-thesis template:
%    - \checkmark comes from amssymb, which ntua-thesis.cls loads in
%      BOTH font branches (cls lines 687 and 692), so it is present
%      whether you compile with "modern" or "classic".
%    - \textcolor comes from xcolor, pulled in by the unconditional
%      \usepackage{tikz} at cls line 1238.
%    - Rules use plain \hline instead of booktabs.
%    - The 3-line header is built by hand instead of with multirow.
%
%  Use related_work_table.tex instead if you install pifont, booktabs
%  and multirow -- it looks better (proper booktabs rule weights).
%
%  USAGE
%    % =====================================================================
%  related_work_table_nodeps.tex
%  Landscape of accelerators for structured sparse GEMV.
%
%  ZERO-DEPENDENCY VARIANT of related_work_table.tex.
%
%  This version needs NO new packages in the ntua-thesis template:
%    - \checkmark comes from amssymb, which ntua-thesis.cls loads in
%      BOTH font branches (cls lines 687 and 692), so it is present
%      whether you compile with "modern" or "classic".
%    - \textcolor comes from xcolor, pulled in by the unconditional
%      \usepackage{tikz} at cls line 1238.
%    - Rules use plain \hline instead of booktabs.
%    - The 3-line header is built by hand instead of with multirow.
%
%  Use related_work_table.tex instead if you install pifont, booktabs
%  and multirow -- it looks better (proper booktabs rule weights).
%
%  USAGE
%    % =====================================================================
%  related_work_table_nodeps.tex
%  Landscape of accelerators for structured sparse GEMV.
%
%  ZERO-DEPENDENCY VARIANT of related_work_table.tex.
%
%  This version needs NO new packages in the ntua-thesis template:
%    - \checkmark comes from amssymb, which ntua-thesis.cls loads in
%      BOTH font branches (cls lines 687 and 692), so it is present
%      whether you compile with "modern" or "classic".
%    - \textcolor comes from xcolor, pulled in by the unconditional
%      \usepackage{tikz} at cls line 1238.
%    - Rules use plain \hline instead of booktabs.
%    - The 3-line header is built by hand instead of with multirow.
%
%  Use related_work_table.tex instead if you install pifont, booktabs
%  and multirow -- it looks better (proper booktabs rule weights).
%
%  USAGE
%    \input{tables/related_work_table_nodeps}
% =====================================================================

\providecommand{\cmark}{\textcolor{green!60!black}{$\checkmark$}}
\providecommand{\xmark}{\textcolor{red}{$\times$}}

\begin{table}[t]
\centering
\small
\setlength{\tabcolsep}{4pt}
\begin{tabular}{lccccccc}
\hline
             & Struct. & Runtime        & Group        & HBM     & No      & Real Si.\ + & Iso-Platform \\
             & $N{:}M$ & Ratio          & $M\!\ge\!32$ & Multi-  & Reduct. & Meas.       & Dense        \\
Architecture & Weights & Switch         & (93.75\%)    & Channel & Tree    & Power       & Baseline     \\
\hline
Mishra~\cite{mishra2021}        & \cmark & \xmark          & \xmark & \xmark & \cmark & \xmark & \cmark \\
FNM-Trans~\cite{fnmtrans2024}   & \cmark & \cmark\,(lim.)  & \xmark & \xmark & \xmark & \xmark & \xmark \\
Serpens~\cite{serpens2022}      & \xmark & \xmark          & \xmark & \cmark & \xmark & \cmark & \xmark \\
SST~\cite{sst2025}              & \cmark & \cmark          & \xmark & \xmark & \xmark & \xmark & \xmark \\
Transitive Array~\cite{ta2025}  & \xmark & \xmark          & \xmark & \xmark & \xmark & \xmark & \xmark \\
DFX~\cite{dfx2022}              & \xmark & \xmark          & \xmark & \cmark & \xmark & \cmark & \xmark \\
DynaX~\cite{dynax2025}          & \cmark & \cmark\,(lim.)  & \cmark & \xmark & \xmark & \xmark & \xmark \\
FlightVGM~\cite{flightvgm2025}  & \xmark & \cmark          & \xmark & \cmark & \xmark & \cmark & \xmark \\
\hline
\textbf{This work}              & \cmark & \cmark          & \cmark & \cmark & \cmark & \cmark & \cmark \\
\hline
\end{tabular}
\caption{Landscape of accelerators for structured sparse GEMV.
\emph{Struct.\ $N{:}M$ Weights}: a fixed, hardware-friendly sparsity
pattern imposed on the weight tensor.
\emph{Runtime Ratio Switch}: the sparsity ratio changes without
re-synthesis.
\emph{Group $M\!\ge\!32$}: the group size reaches 32, i.e.\ a 93.75\%
sparsity ceiling rather than the 75\% of $M\!=\!4$ designs.
\emph{HBM Multi-Channel}: operands streamed over HBM pseudo-channels.
\emph{No Reduction Tree}: load balance by construction, with no
arbiters, schedulers, crossbars or adder trees.
\emph{Real Si.\ + Meas.\ Power}: deployed on silicon with power
measured in-system rather than modelled or read from a synthesis
report.
\emph{Iso-Platform Dense Baseline}: sparse compared against a dense
counterpart on the same board, flow and clock.}
\label{tab:related-work}
\end{table}

% =====================================================================

\providecommand{\cmark}{\textcolor{green!60!black}{$\checkmark$}}
\providecommand{\xmark}{\textcolor{red}{$\times$}}

\begin{table}[t]
\centering
\small
\setlength{\tabcolsep}{4pt}
\begin{tabular}{lccccccc}
\hline
             & Struct. & Runtime        & Group        & HBM     & No      & Real Si.\ + & Iso-Platform \\
             & $N{:}M$ & Ratio          & $M\!\ge\!32$ & Multi-  & Reduct. & Meas.       & Dense        \\
Architecture & Weights & Switch         & (93.75\%)    & Channel & Tree    & Power       & Baseline     \\
\hline
Mishra~\cite{mishra2021}        & \cmark & \xmark          & \xmark & \xmark & \cmark & \xmark & \cmark \\
FNM-Trans~\cite{fnmtrans2024}   & \cmark & \cmark\,(lim.)  & \xmark & \xmark & \xmark & \xmark & \xmark \\
Serpens~\cite{serpens2022}      & \xmark & \xmark          & \xmark & \cmark & \xmark & \cmark & \xmark \\
SST~\cite{sst2025}              & \cmark & \cmark          & \xmark & \xmark & \xmark & \xmark & \xmark \\
Transitive Array~\cite{ta2025}  & \xmark & \xmark          & \xmark & \xmark & \xmark & \xmark & \xmark \\
DFX~\cite{dfx2022}              & \xmark & \xmark          & \xmark & \cmark & \xmark & \cmark & \xmark \\
DynaX~\cite{dynax2025}          & \cmark & \cmark\,(lim.)  & \cmark & \xmark & \xmark & \xmark & \xmark \\
FlightVGM~\cite{flightvgm2025}  & \xmark & \cmark          & \xmark & \cmark & \xmark & \cmark & \xmark \\
\hline
\textbf{This work}              & \cmark & \cmark          & \cmark & \cmark & \cmark & \cmark & \cmark \\
\hline
\end{tabular}
\caption{Landscape of accelerators for structured sparse GEMV.
\emph{Struct.\ $N{:}M$ Weights}: a fixed, hardware-friendly sparsity
pattern imposed on the weight tensor.
\emph{Runtime Ratio Switch}: the sparsity ratio changes without
re-synthesis.
\emph{Group $M\!\ge\!32$}: the group size reaches 32, i.e.\ a 93.75\%
sparsity ceiling rather than the 75\% of $M\!=\!4$ designs.
\emph{HBM Multi-Channel}: operands streamed over HBM pseudo-channels.
\emph{No Reduction Tree}: load balance by construction, with no
arbiters, schedulers, crossbars or adder trees.
\emph{Real Si.\ + Meas.\ Power}: deployed on silicon with power
measured in-system rather than modelled or read from a synthesis
report.
\emph{Iso-Platform Dense Baseline}: sparse compared against a dense
counterpart on the same board, flow and clock.}
\label{tab:related-work}
\end{table}

% =====================================================================

\providecommand{\cmark}{\textcolor{green!60!black}{$\checkmark$}}
\providecommand{\xmark}{\textcolor{red}{$\times$}}

\begin{table}[t]
\centering
\small
\setlength{\tabcolsep}{4pt}
\begin{tabular}{lccccccc}
\hline
             & Struct. & Runtime        & Group        & HBM     & No      & Real Si.\ + & Iso-Platform \\
             & $N{:}M$ & Ratio          & $M\!\ge\!32$ & Multi-  & Reduct. & Meas.       & Dense        \\
Architecture & Weights & Switch         & (93.75\%)    & Channel & Tree    & Power       & Baseline     \\
\hline
Mishra~\cite{mishra2021}        & \cmark & \xmark          & \xmark & \xmark & \cmark & \xmark & \cmark \\
FNM-Trans~\cite{fnmtrans2024}   & \cmark & \cmark\,(lim.)  & \xmark & \xmark & \xmark & \xmark & \xmark \\
Serpens~\cite{serpens2022}      & \xmark & \xmark          & \xmark & \cmark & \xmark & \cmark & \xmark \\
SST~\cite{sst2025}              & \cmark & \cmark          & \xmark & \xmark & \xmark & \xmark & \xmark \\
Transitive Array~\cite{ta2025}  & \xmark & \xmark          & \xmark & \xmark & \xmark & \xmark & \xmark \\
DFX~\cite{dfx2022}              & \xmark & \xmark          & \xmark & \cmark & \xmark & \cmark & \xmark \\
DynaX~\cite{dynax2025}          & \cmark & \cmark\,(lim.)  & \cmark & \xmark & \xmark & \xmark & \xmark \\
FlightVGM~\cite{flightvgm2025}  & \xmark & \cmark          & \xmark & \cmark & \xmark & \cmark & \xmark \\
\hline
\textbf{This work}              & \cmark & \cmark          & \cmark & \cmark & \cmark & \cmark & \cmark \\
\hline
\end{tabular}
\caption{Landscape of accelerators for structured sparse GEMV.
\emph{Struct.\ $N{:}M$ Weights}: a fixed, hardware-friendly sparsity
pattern imposed on the weight tensor.
\emph{Runtime Ratio Switch}: the sparsity ratio changes without
re-synthesis.
\emph{Group $M\!\ge\!32$}: the group size reaches 32, i.e.\ a 93.75\%
sparsity ceiling rather than the 75\% of $M\!=\!4$ designs.
\emph{HBM Multi-Channel}: operands streamed over HBM pseudo-channels.
\emph{No Reduction Tree}: load balance by construction, with no
arbiters, schedulers, crossbars or adder trees.
\emph{Real Si.\ + Meas.\ Power}: deployed on silicon with power
measured in-system rather than modelled or read from a synthesis
report.
\emph{Iso-Platform Dense Baseline}: sparse compared against a dense
counterpart on the same board, flow and clock.}
\label{tab:related-work}
\end{table}

% =====================================================================

\providecommand{\cmark}{\textcolor{green!60!black}{$\checkmark$}}
\providecommand{\xmark}{\textcolor{red}{$\times$}}

\begin{table}[t]
\centering
\small
\setlength{\tabcolsep}{4pt}
\begin{tabular}{lccccccc}
\hline
             & Struct. & Runtime        & Group        & HBM     & No      & Real Si.\ + & Iso-Platform \\
             & $N{:}M$ & Ratio          & $M\!\ge\!32$ & Multi-  & Reduct. & Meas.       & Dense        \\
Architecture & Weights & Switch         & (93.75\%)    & Channel & Tree    & Power       & Baseline     \\
\hline
Mishra~\cite{mishra2021}        & \cmark & \xmark          & \xmark & \xmark & \cmark & \xmark & \cmark \\
FNM-Trans~\cite{fnmtrans2024}   & \cmark & \cmark\,(lim.)  & \xmark & \xmark & \xmark & \xmark & \xmark \\
Serpens~\cite{serpens2022}      & \xmark & \xmark          & \xmark & \cmark & \xmark & \cmark & \xmark \\
SST~\cite{sst2025}              & \cmark & \cmark          & \xmark & \xmark & \xmark & \xmark & \xmark \\
Transitive Array~\cite{ta2025}  & \xmark & \xmark          & \xmark & \xmark & \xmark & \xmark & \xmark \\
DFX~\cite{dfx2022}              & \xmark & \xmark          & \xmark & \cmark & \xmark & \cmark & \xmark \\
DynaX~\cite{dynax2025}          & \cmark & \cmark\,(lim.)  & \cmark & \xmark & \xmark & \xmark & \xmark \\
FlightVGM~\cite{flightvgm2025}  & \xmark & \cmark          & \xmark & \cmark & \xmark & \cmark & \xmark \\
\hline
\textbf{This work}              & \cmark & \cmark          & \cmark & \cmark & \cmark & \cmark & \cmark \\
\hline
\end{tabular}
\caption{Landscape of accelerators for structured sparse GEMV.
\emph{Struct.\ $N{:}M$ Weights}: a fixed, hardware-friendly sparsity
pattern imposed on the weight tensor.
\emph{Runtime Ratio Switch}: the sparsity ratio changes without
re-synthesis.
\emph{Group $M\!\ge\!32$}: the group size reaches 32, i.e.\ a 93.75\%
sparsity ceiling rather than the 75\% of $M\!=\!4$ designs.
\emph{HBM Multi-Channel}: operands streamed over HBM pseudo-channels.
\emph{No Reduction Tree}: load balance by construction, with no
arbiters, schedulers, crossbars or adder trees.
\emph{Real Si.\ + Meas.\ Power}: deployed on silicon with power
measured in-system rather than modelled or read from a synthesis
report.
\emph{Iso-Platform Dense Baseline}: sparse compared against a dense
counterpart on the same board, flow and clock.}
\label{tab:related-work}
\end{table}
